\documentclass[11pt]{article}
\usepackage{amsmath, amssymb, amsthm}
\usepackage{geometry}
\usepackage{hyperref}
\usepackage{booktabs}
\usepackage{xcolor}

\geometry{margin=1.1in}

\newcommand{\eps}{\varepsilon}
\newcommand{\dt}{\Delta t}
\newcommand{\dx}{\Delta x}
\newcommand{\RR}{\mathbb{R}}
\newcommand{\EE}{\mathbb{E}}
\newcommand{\Ht}{H_t}
\newcommand{\HT}{H_T}

\theoremstyle{remark}
\newtheorem{remark}{Remark}

\title{Dynamic Sinkhorn for the Schr\"odinger Bridge:\\
       Free-Space vs.\ Neumann Heat Kernels}
\author{Optimal Transport / Schr\"odinger Bridge Project}
\date{}

\begin{document}
\maketitle

\begin{abstract}
The Sinkhorn algorithm for the dynamic Schr\"odinger bridge (SB) requires applying a heat
semigroup at each iteration.
The choice of heat kernel determines what reference stochastic process the SB is built
against, and how the kernel is computed in practice.
This note documents two implementations:
(1) the Neumann kernel on $[0,1]$ (reflecting Brownian motion, computed via DCT), and
(2) the free-space kernel on $\RR$ (standard Brownian motion, computed via FFT
convolution).
We describe the mathematics, the discrete implementations, the normalisation
conventions, and when the two kernels agree or differ.
\end{abstract}

%------------------------------------------------------------------------------
\section{The Schr\"odinger Bridge Problem}
%------------------------------------------------------------------------------

Given two probability densities $\rho_0, \rho_1$ on a domain $\Omega$
(either $[0,1]$ or $\RR$) and diffusivity $\eps > 0$, the Schr\"odinger bridge
problem asks for the probability path measure $P^*$ that minimises the KL
divergence to the Wiener measure $W^\eps$ (Brownian motion with diffusivity $\eps$)
subject to having marginals $\rho_0$ at $t=0$ and $\rho_1$ at $t=T$:
\begin{equation}
  P^* = \arg\min_{P:\, \rho_P(0)=\rho_0,\, \rho_P(T)=\rho_1} \mathrm{KL}(P \| W^\eps).
\end{equation}

The hydrodynamic (PDE) form of the same problem is:
\begin{equation}
  \min_{\rho, m} \;\iint_{\Omega} \frac{|m|^2}{2\rho} + \eps^2 \rho \log \rho \; dx\, dt
  \quad \text{s.t.} \quad \partial_t \rho + \partial_x m = \eps \Delta\rho, \quad
  \rho(0,\cdot)=\rho_0, \quad \rho(T,\cdot)=\rho_1.
\end{equation}
(The $\eps^2 \rho \log \rho$ term is the entropic regularisation; for $\eps \to 0$ this
vanishes and we recover the Benamou--Brenier OT problem.)

%------------------------------------------------------------------------------
\section{Hopf--Cole Factorisation and the Schr\"odinger System}
%------------------------------------------------------------------------------

The optimal density factorises as
\begin{equation}
  \rho^*(t, x) = \varphi(t,x) \cdot \psi(t,x),
\end{equation}
where $\varphi$ satisfies the \textit{forward} heat equation and $\psi$ the
\textit{backward} heat equation:
\begin{equation}
  \partial_t \varphi = \eps \Delta \varphi, \qquad
  -\partial_t \psi  = \eps \Delta \psi.
\end{equation}
The pair $(\varphi, \psi)$ must satisfy the \textbf{Schr\"odinger system}:
\begin{equation}
  \varphi(0,x)\, \psi(0,x) = \rho_0(x), \qquad
  \varphi(T,x)\, \psi(T,x) = \rho_1(x).
\end{equation}

\textbf{Momentum.}
The optimal momentum density is recovered from the Hopf--Cole factors as
\begin{equation}
  m^*(t,x) = 2\eps\, \varphi(t,x)\, \partial_x \psi(t,x).
\end{equation}
In the code the spatial gradient is computed at staggered positions $x_{j+1/2}$:
\begin{equation}
  m^*(t, x_{j+\frac{1}{2}})
  = 2\eps\, \frac{\varphi_j(t) + \varphi_{j+1}(t)}{2}
             \cdot \frac{\psi_{j+1}(t) - \psi_j(t)}{\dx}.
\end{equation}

\textbf{Trajectory after convergence.}
Given converged boundary values $\varphi(0,\cdot)$ and $\psi(T,\cdot)$, the full
trajectory is recovered by propagating each forward/backward in time:
\begin{equation}
  \varphi(t,\cdot) = H_t[\varphi(0,\cdot)], \qquad
  \psi(t,\cdot) = H_{T-t}[\psi(T,\cdot)],
\end{equation}
using the heat semigroup $H_t$ of the chosen kernel.

%------------------------------------------------------------------------------
\section{Sinkhorn Iterations}
%------------------------------------------------------------------------------

Let $H_T$ denote the heat semigroup at the full time $T$.
Starting from $\psi^{(0)}(0,\cdot) = \mathbf{1}$, each Sinkhorn iteration is:
\begin{align}
  \varphi^{(n+1)}(0) &\leftarrow \rho_0 / \psi^{(n)}(0)       \label{eq:sink1} \\
  \varphi^{(n+1)}(T) &= H_T[\varphi^{(n+1)}(0)]               \label{eq:sink2} \\
  \psi^{(n+1)}(T)    &\leftarrow \rho_1 / \varphi^{(n+1)}(T)  \label{eq:sink3} \\
  \psi^{(n+1)}(0)    &= H_T[\psi^{(n+1)}(T)]                  \label{eq:sink4}
\end{align}
where division is pointwise.
Convergence is measured as the $L^2$ error on the left marginal:
\begin{equation}
  \text{err}^{(n)} = \|\varphi^{(n)}(0,\cdot)\, \psi^{(n)}(0,\cdot) - \rho_0\|_{L^2}.
\end{equation}
Convergence is guaranteed for any $\eps > 0$ and is linear with rate determined by the
``Hilbert metric'' contraction of $H_T$.
The rate worsens as $\eps \to 0$ (less regularisation), making Sinkhorn impractical for
pure OT ($\eps = 0$).

\subsection{One-Shot vs.\ Step-by-Step Propagation}

There are two ways to apply $H_T$ numerically:
\begin{itemize}
  \item \textbf{One-shot}: apply $H_T$ directly as a single operation (one convolution
        with $K_T$, or one DCT multiplication by $e^{-\eps\lambda_k T}$).
  \item \textbf{Step-by-step}: compose $n_t$ applications of $H_{\dt}$ (smaller
        convolutions or DCT multiplications).
\end{itemize}

For the Neumann kernel both are equivalent (both exact up to floating point), since the
DCT diagonalises the semigroup for all $t$.

For the free-space kernel on a bounded grid, only one-shot is correct.
$\varphi(0)$ and $\psi(T)$ are narrow (they equal the marginal divided by a slowly
varying function), so a single zero-padded convolution with $K_T$ is accurate.
With step-by-step propagation, after $k$ steps the function has spread wide and is
truncated by the grid boundary at each subsequent step; the error compounds over
$n_t$ steps and can be $O(1)$ for large $\eps$.

The current implementation uses one-shot propagation inside the Sinkhorn loop,
and step-by-step propagation (with the smaller kernel $H_{T-t_k}$ applied once per
output time) during trajectory recovery.

%------------------------------------------------------------------------------
\section{Kernel 1: Neumann on $[0,1]$ (DCT)}
%------------------------------------------------------------------------------

\subsection{Reference Process}

The reference process is \textbf{reflected Brownian motion on $[0,1]$}: standard
Brownian motion with diffusivity $\eps$ that reflects elastically at $x = 0$ and $x = 1$.
The associated SB is the optimal coupling of $\rho_0$ and $\rho_1$ relative to this
reflected process.

\subsection{Heat Semigroup}

The heat equation $\partial_t u = \eps \Delta u$ with homogeneous Neumann BCs
($\partial_x u|_{x=0} = \partial_x u|_{x=1} = 0$) is diagonalised by the cosine functions:
\begin{equation}
  e_k(x) = \cos(k \pi x), \qquad k = 0, 1, 2, \ldots, n_x - 1.
\end{equation}
The eigenvalue of $-\Delta$ for mode $k$ is $\lambda_k = (k\pi)^2$ (continuous).
The heat semigroup acts by
\begin{equation}
  \bigl(H_t \varphi\bigr)(x)
  = \sum_{k=0}^{\infty} e^{-\eps \lambda_k t}\, \hat\varphi_k\, e_k(x),
\end{equation}
where $\hat\varphi_k = \langle \varphi, e_k \rangle$ are the cosine (DCT) coefficients.

\subsection{Discrete Implementation}

On an $n_x$-point cell-centre grid with $\dx = 1/n_x$, the DCT-II diagonalises the
\textit{discrete} Neumann Laplacian.
The discrete eigenvalues are
\begin{equation}
  \lambda_k^{\dx} = \frac{2 - 2\cos(k\pi\dx)}{\dx^2},
  \qquad k = 0, 1, \ldots, n_x - 1.
\end{equation}
For small $k\dx$: $\lambda_k^{\dx} \approx (k\pi)^2 = \lambda_k$, so the discrete and
continuous eigenvalues agree at low frequencies and diverge near the Nyquist mode
$k = n_x - 1$.

The semigroup application reduces to three operations:
\begin{equation}
  H_t[\varphi] = \mathrm{IDCT}\!\left(e^{-\eps\,\lambda^{\dx}\, t} \odot \mathrm{DCT}(\varphi)\right),
\end{equation}
where $\odot$ is elementwise multiplication.
The decay vector $e^{-\eps\,\lambda^{\dx}\, T}$ (for the Sinkhorn loop) is precomputed once.
For trajectory recovery at time $t$, the decay $e^{-\eps\,\lambda^{\dx}\,t}$ is recomputed
from the eigenvalues, which are also precomputed.

\subsection{Marginal Normalisation}

For the Neumann kernel the marginals $\rho_0$, $\rho_1$ must be normalised as
\textbf{discrete probability masses} summing to 1:
\begin{equation}
  \rho_0^{(i)} = \frac{f(x_i)}{\sum_j f(x_j)}, \qquad \sum_i \rho_0^{(i)} = 1.
\end{equation}
The Neumann SB on $[0,1]$ is a \textit{coupling} of these mass vectors.
The steady state of $H_t$ as $t \to \infty$ is the uniform distribution (equal mass at
each grid point), which is the correct equilibrium for reflected BM.

\subsection{When to Use}

\begin{itemize}
  \item Marginals are well-concentrated in the interior of $[0,1]$
        (i.e.\ $\sigma \ll \min(\mu, 1-\mu)$), so the Neumann and free-space kernels
        agree to high accuracy.
  \item The problem is naturally set on a bounded domain with no-flux boundaries.
  \item Numerical efficiency: DCT is $O(n_x \log n_x)$, no padding overhead.
\end{itemize}

%------------------------------------------------------------------------------
\section{Kernel 2: Free Space on $\RR$ (Gaussian Convolution)}
%------------------------------------------------------------------------------

\subsection{Reference Process}

The reference process is \textbf{standard Brownian motion on $\RR$} with diffusivity
$\eps$.
The associated SB is the optimal coupling of $\rho_0$ and $\rho_1$ on all of $\RR$,
without boundary effects.
This is the mathematically natural setting for the Gaussian SB test case, where the
analytical solution is known explicitly.

\subsection{Heat Semigroup}

The heat kernel on $\RR$ is the Gaussian:
\begin{equation}
  K_t(x, y) = \frac{1}{\sqrt{4\pi\eps t}} \exp\!\left(-\frac{(x-y)^2}{4\eps t}\right).
\end{equation}
The semigroup acts by convolution:
\begin{equation}
  (H_t \varphi)(x) = \int_\RR K_t(x-y)\, \varphi(y)\, dy
                   = (K_t * \varphi)(x).
\end{equation}
Note that $K_t$ is a Gaussian with variance $\sigma^2 = 2\eps t$:
$K_t \sim \mathcal{N}(0, 2\eps t)$.

\subsection{Discrete Implementation}

On a finite grid $\{x_i\}_{i=1}^{n_x} \subset [0,1]$ the convolution is approximated by
zero-padded FFT convolution.
The zero-padding is essential to avoid wrap-around aliasing from the cyclic FFT.
The padding width is chosen so that the Gaussian tail is negligible at the boundary:
\begin{equation}
  n_{\mathrm{pad}} = \left\lceil \frac{6\sqrt{2\eps T}}{\dx} \right\rceil
  \qquad (\text{6 standard deviations}).
\end{equation}
The padded length is rounded up to the next power of 2 for FFT efficiency:
$n_{\mathrm{fft}} = 2^{\lceil \log_2(n_x + 2 n_{\mathrm{pad}}) \rceil}$.

The Gaussian kernel is discretised on the padded grid, normalised so that
$\sum_j K(j\dx) = 1$ (discrete probability), and FFT'd once per precomputation:
\begin{equation}
  H_T[\varphi] = \mathrm{Re}\!\left(
    \mathcal{F}^{-1}\!\left[\mathcal{F}[\varphi_{\mathrm{pad}}] \odot \hat K_T\right]
  \right)\!\big|_{[1:n_x]}.
\end{equation}
The output is cropped back to the original $n_x$ grid after convolution.

For trajectory recovery, the kernel is recomputed at each time $t$ (not precomputed).

\subsection{Marginal Normalisation}

For the free-space kernel the marginals $\rho_0$, $\rho_1$ must be supplied as
\textbf{PDF values} (not normalised to sum to 1):
\begin{equation}
  \rho_0^{(i)} = f(x_i), \qquad \sum_i \rho_0^{(i)}\,\dx \approx 1.
\end{equation}
This is because the free-space SB is a coupling of continuous distributions on $\RR$,
not of discrete masses on $[0,1]$.
The steady state of $K_t$ as $t \to \infty$ is the zero function (mass escapes to
$\pm\infty$), not a uniform distribution.

In the code, the flag \texttt{cfg.use\_pdf\_marginals = true} selects this convention;
the default (\texttt{false}) uses the mass-normalised convention for the Neumann kernel.

\subsection{When to Use}

\begin{itemize}
  \item Marginals are defined on $\RR$ (e.g.\ Gaussians), and the analytical SB
        solution is known for comparison.
  \item Marginals extend to or are concentrated near the boundary of $[0,1]$, making
        the Neumann reflecting BCs physically incorrect.
  \item Validation: the free-space SB for Gaussian marginals has a closed-form solution
        (linear interpolation of means and standard deviations), enabling exact error measurement.
\end{itemize}

%------------------------------------------------------------------------------
\section{Comparison}
%------------------------------------------------------------------------------

\begin{center}
\renewcommand{\arraystretch}{1.4}
\begin{tabular}{lll}
\toprule
                        & \textbf{Neumann (DCT)}     & \textbf{Free Space (FFT)} \\
\midrule
Reference process       & Reflected BM on $[0,1]$    & Standard BM on $\RR$      \\
Spatial domain          & $[0,1]$                    & $\RR$ (truncated to $[0,1]$) \\
BCs at walls            & Reflecting (zero flux)     & None (absorbing / escaping) \\
Kernel formula          & $e^{-\eps\lambda_k t}$ in DCT basis & Gaussian $K_t(x) = \mathcal{N}(0, 2\eps t)$ \\
Implementation          & DCT + pointwise multiply   & Zero-padded FFT convolution \\
Precomputation          & Decay vector $e^{-\eps\lambda^{\dx} T}$ & Kernel FFT $\hat K_T$ \\
Cost per Sinkhorn iter  & $O(n_x \log n_x)$          & $O(n_{\mathrm{fft}} \log n_{\mathrm{fft}})$, $n_{\mathrm{fft}} \gg n_x$ \\
Marginal normalisation  & Mass (sum = 1)             & PDF (sum$\cdot\dx \approx 1$) \\
Steady state of $H_t$   & Uniform mass               & Zero (mass escapes) \\
Analytical test case    & No (known only for $\eps \to 0$) & Yes (Gaussian SB) \\
\bottomrule
\end{tabular}
\end{center}

\subsection{When Do the Two Kernels Agree?}

The Neumann kernel on $[0,1]$ and the free-space kernel on $\RR$ agree when the
marginals are concentrated far from the boundary.
More precisely, by the \textbf{method of images}, the Neumann heat kernel is the
free-space heat kernel plus image charges reflected at $x = 0$ and $x = 1$:
\begin{equation}
  K_t^{\mathrm{Neumann}}(x, y)
  = \sum_{n=-\infty}^{\infty}
    \bigl[ K_t(x, y + 2n) + K_t(x, -y + 2n) \bigr].
\end{equation}
The correction terms decay like $e^{-1/(2\eps t)}$ when $y$ is bounded away from 0 and 1.
For a Gaussian marginal $\mathcal{N}(\mu, \sigma^2)$ with $\sigma \ll \min(\mu, 1-\mu)$,
the image charges contribute exponentially small corrections, and the two kernels are
numerically indistinguishable.
This is the regime of our Gaussian test cases (e.g.\ $\mu = 1/3$ or $2/3$, $\sigma = 0.05$).

\subsection{When Do They Differ?}

They differ when:
\begin{itemize}
  \item Marginals have significant mass near the boundary (large $\sigma$, $\mu$ near 0 or 1).
  \item $\eps$ is large enough that the heat kernel has spread to the walls
        over time $T$: roughly $\sqrt{2\eps T} \gtrsim \min(\mu, 1-\mu)$.
  \item The problem is genuinely set on $[0,1]$ with physical no-flux walls
        (Neumann is correct), or genuinely on $\RR$ (free-space is correct).
\end{itemize}

%------------------------------------------------------------------------------
\section{Implementation Files}
%------------------------------------------------------------------------------

\begin{itemize}
  \item \texttt{shared/1d/pipelines/sinkhorn\_hopf\_cole.m} ---
        main Sinkhorn loop, kernel-agnostic via \texttt{cfg.precomp\_heat}.
  \item \texttt{shared/1d/utils/precomp\_heat\_neumann.m} ---
        builds the Neumann (DCT) kernel handle.
  \item \texttt{shared/1d/utils/precomp\_heat\_free\_space.m} ---
        builds the free-space (Gaussian FFT) kernel handle.
  \item \texttt{discretize\_first/1d/experiments/test\_sb\_gaussian\_sinkhorn.m} ---
        comparison of ADMM and Sinkhorn (both kernels) on the Gaussian SB test case.
\end{itemize}

To switch kernels, change \texttt{cfg.precomp\_heat}:
\begin{verbatim}
  cfg.precomp_heat = @precomp_heat_neumann;      % reflecting BM on [0,1]
  cfg.precomp_heat = @precomp_heat_free_space;   % BM on R
\end{verbatim}
For the free-space kernel also set \texttt{cfg.use\_pdf\_marginals = true} and
use \texttt{problem.rho0\_pdf}/\texttt{rho1\_pdf} rather than the mass-normalised
\texttt{rho0}/\texttt{rho1}.

%------------------------------------------------------------------------------
\section{Current Challenges}
%------------------------------------------------------------------------------

\begin{itemize}
  \item \textbf{Trajectory recovery cost.}
        The free-space trajectory recovery recomputes a new Gaussian kernel for each of
        the $n_t + 1$ output times.
        For large $n_t$ this is non-trivial; the kernels could be precomputed once (as
        is already done for the Neumann kernel).

  \item \textbf{Step-by-step Neumann for intermediate times.}
        During trajectory recovery with the Neumann kernel, the code calls
        \texttt{heat.apply\_time($\varphi_0$, $t_k$)} which recomputes
        $e^{-\eps\lambda^{\dx} t_k}$ at each output time.
        For large $n_t$ this could be replaced by stepping forward one $\dt$ at a time
        using the precomputed \texttt{decay\_dt}, which only requires one multiply per step.

  \item \textbf{Extension to 2D.}
        The Neumann kernel extends naturally to 2D via a 2D DCT (diagonalising the 2D
        Neumann Laplacian with eigenvalues $\lambda_{kx} + \lambda_{ky}$).
        The free-space kernel extends via 2D Gaussian convolution with a 2D FFT.
        Neither is yet implemented; the current 2D solver uses ADMM only.
\end{itemize}

\end{document}
